\subsection{Arbitrary finite Frobenius data coexist with real-multiplication theta avoidance}
\label{sec:kt-rm-frobenius-avoidance}

The following is a strengthening of the \emph{carried method} in
AN\_GRH~A37, rather than a new claim about its published moduli
inputs.  A37 constructs avoidance with one specially selected
Frobenius class.  The additional argument below permits any
prescribed finite symplectic Frobenius class, including at the
primes to be avoided.  It identifies a limit of arguments using
finite torsion data; it does not address an exact Weil polynomial
or a fixed ground-field size.

Put $\mathcal O=\mathbf Z[T]/(T^2+T-1)=\mathcal O_5$.
All polarizations below are the canonical principal polarizations,
and the specified $\mathcal O$ action is fixed by Rosati.

\begin{theorem}[KT-RM-FROB-AVOID]
Let $m\geq1$ be odd, let $S$ be a finite set of odd primes, and
let
\[
                    \gamma\in\operatorname{SL}_2(\mathcal O/m\mathcal O).
                                                               \tag{RF.1}
\]
In every prescribed nonempty open of the ordinary genus-two
Hilbert--Blumenthal locus in characteristic two there is a curve
$C/\mathbf F_Q$,
with its specified integral Rosati-$\mathcal O$ action defined
over $\mathbf F_Q$, such that:
\begin{enumerate}
\item $Q\equiv1\pmod m$, and the action of $Q$-power Frobenius
      on the rank-two $\mathcal O/m$ module
      $J=\operatorname{Pic}^0(C^{(1)})[m]$ is conjugate to $\gamma$
      in an $\mathcal O$-linear symplectic frame;
\item the actual Raynaud theta support satisfies
\[
 \Theta_C(\overline{\mathbf F}_Q)\cap
 \bigl(J_{C^{(1)}}(\overline{\mathbf F}_Q)\{S\}\setminus\{0\}\bigr)
                          =\varnothing.                       \tag{RF.2}
\]
\end{enumerate}
The notation $\{S\}$ permits every exponent and every product
of primes in $S$.  No disjointness between $S$ and the prime
divisors of $m$ is assumed.  No bound for $Q$ is asserted.
\end{theorem}

\begin{proof}
We use two precise carried inputs.  AN\_GRH~A37, Steps~1--4,
gives an ordinary genus-two universal family on a geometrically
irreducible Hilbert--Blumenthal level scheme, full joint monodromy
$\operatorname{SL}_2(\mathcal O/n)$ at every odd level $n$, and
a nonempty open avoiding nonzero theta $n$-torsion for each fixed
$n$.  Its inverse-different pairing convention includes the
ramified prime five.  AN\_GRH~A19 gives the following actual
geometric statement: if $\Phi^d$ fixes $J[\ell]$ for every
$\ell\in S$, then every $S$-primary point on a translated genus-two
Abel curve defined over $\mathbf F_Q$ lies in
$J(\mathbf F_{Q^d})$.  It does not require semisimple Frobenius.

\paragraph{1. A lift of any finite class with bounded fixed torsion.}
At each $\ell\in S$, prescribe the reduction of $\gamma$ modulo
$\ell^{v_\ell(m)}$ when $\ell\mid m$; when $\ell\nmid m$,
prescribe the identity modulo $\ell$.  Let $D\geq1$ be a common
multiple of the orders of the resulting actions on the underlying
rank-four $\mathbf F_\ell$ modules.  Thus every lift $G_\ell$
of these prescribed classes satisfies $G_\ell^D\equiv1\pmod\ell$.

There is such a lift
\[
 G_\ell\in\operatorname{SL}_2(\mathcal O\otimes\mathbf Z_\ell)
 \quad\hbox{with}\quad
             \det_{\mathbf Z_\ell}(G_\ell^D-1)\ne0.          \tag{RF.3}
\]
Here the determinant is on the underlying rank-four module.
We give the existence argument explicitly.  Reduction maps for
$\operatorname{SL}_2$ over the local factors of
$\mathcal O\otimes\mathbf Z_\ell$ are surjective: lift a matrix
with determinant one and correct an entry using a unit cofactor,
then pass through the successive quotients.  This works at five
as well as at unramified primes.  Fix one lift $G_0$.

The regular function
$g\mapsto\det_{\mathbf Q_\ell}(g^D-1)$ on
$\operatorname{Res}_{F\otimes\mathbf Q_\ell/\mathbf Q_\ell}
\operatorname{SL}_2$, where $F=\mathbf Q(\sqrt5)$, is not zero.
For example $g=\operatorname{diag}(a,a^{-1})$, with a rational
$\ell$-adic unit $a$ of infinite order, has nonzero value.
Every congruence coset of this group is Zariski dense.  An
elementary verification uses the big cell
\[
 G_0\begin{pmatrix}
       1+x&y\\ z&(1+yz)/(1+x)
     \end{pmatrix},\qquad
 x,y,z\in\ell^r(\mathcal O\otimes\mathbf Z_\ell),             \tag{RF.4}
\]
for sufficiently large $r$.  Over $\mathbf Q_\ell$ this is an
open dense cell in an irreducible algebraic group; after choosing
a $\mathbf Z_\ell$ basis of $\mathcal O\otimes\mathbf Z_\ell$,
a nonzero numerator polynomial cannot vanish on the entire
product of six open $\ell$-adic balls.  The last assertion follows
by induction on the number of variables from the finiteness of
the roots of a nonzero one-variable polynomial.  Hence the
prescribed coset contains a point satisfying (RF.3).

Set
\[
 a_\ell=v_\ell\det_{\mathbf Z_\ell}(G_\ell^D-1),
 \qquad N=\prod_{\ell\in S}\ell^{a_\ell}.                  \tag{RF.5}
\]
These are finite integers.  Choose an odd integer $L$ divisible
by $m$ and by every $\ell^{a_\ell+1}$, and prescribe a class
$\delta\in\operatorname{SL}_2(\mathcal O/L)$ which reduces to
$\gamma$ modulo $m$ and to $G_\ell$ modulo
$\ell^{a_\ell+1}$ for every $\ell\in S$.  If
$v_\ell(m)>a_\ell+1$, retain the already prescribed higher
precision of $G_\ell$.  The Chinese remainder theorem gives
$\delta$; outside $S$, retain the prescribed components of
$\gamma$.

For any integral matrix $\Phi$ congruent to $\delta$ at this
level, polynomial continuity of powers and determinants gives
\[
 v_\ell\det(\Phi^D-1)=a_\ell\quad(\ell\in S).              \tag{RF.6}
\]
On a prime-to-two Tate module, the kernel of $\Phi^D-1$ on
$T_\ell J\otimes\mathbf Q_\ell/\mathbf Z_\ell$ is therefore
killed by $\ell^{a_\ell}$.  For completeness, its Smith normal
form has diagonal entries with total valuation $a_\ell$;
each diagonal valuation is at most their sum.  Equivalently,
the adjugate identity gives this annihilator directly.

\paragraph{2. Specialize inside the finite theta-avoidance open.}
Choose an auxiliary odd level $h\geq3$ prime to $L$.
Apply the generic avoidance in A37 to $N$, inside the prescribed
nonempty ordinary open, and then take a nonempty smooth affine
open $U$ on which the actual theta has no nonzero $N$-torsion.
After descending all data to a finite field
$k_0=\mathbf F_{q_0}$, enlarge it so that
$q_0\equiv1\pmod{hL}$.  This does not affect geometric
irreducibility or the theta-avoidance property.

The full level-$L$ torsor $V\to U$ is geometrically integral
with constant group $\operatorname{SL}_2(\mathcal O/L)$, by
the full joint monodromy input in A37.  Twist it by $\delta$
exactly as in A37, Step~5.  The affine descent datum is effective,
and the twist is geometrically integral.  Lang--Weil supplies a
point over $\mathbf F_{q_0^r}$ for every sufficiently large $r$.
Choose $r\equiv1\pmod{\operatorname{ord}(\delta)}$.
The resulting point of $U$ has arithmetic Frobenius class
$\delta$ in its level-$L$ torsor.  Let $C/\mathbf F_Q$ be
the universal curve at this point, with $Q=q_0^r$.
The actual Raynaud line bundle, its theta support, and the
specified real multiplication are defined over this exact field.
Thus item~1 holds, and $\Phi$ satisfies (RF.6).

\paragraph{3. Pass from a finite level to every supported order.}
Since $\Phi^D$ is the identity on each $J[\ell]$ for
$\ell\in S$, A19 puts every geometric $S$-primary theta point
in $J(\mathbf F_{Q^D})$.  By (RF.6) and the annihilator
calculation after it, that rational $S$-primary group is killed
by $N$.  The specialization was chosen to avoid all nonzero
$N$-torsion on theta, so no such point exists.  This proves
(RF.2).  If $S$ is empty, only the torsor specialization is
needed.  All constructions started inside the specified open,
which proves the final locus assertion.
\end{proof}

\begin{corollary}[KT-RM-FROB-NONDETERMINATION]
For every $b\geq1$ there are two ordinary genus-two curves over
finite fields of characteristic two, each with a specified
integral Rosati-$\mathcal O_5$ action, whose Frobenius actions
on $4159^b$-torsion are isomorphic as symplectic
$\mathcal O_5/4159^b$ modules, such that:
\begin{enumerate}
\item the first curve has an actual odd cyclic cover of degree
      $4159$ with primitive two-rank zero;
\item the second curve has no such cover of degree $4159^a$
      for \emph{any} $a\geq1$.
\end{enumerate}
Both field sizes can be required to be $1$ modulo $4159^b$.
No equality of the two field sizes or of their Weil polynomials
is asserted.  The common Frobenius class is chosen from the
actual positive curve; this corollary does not assert a positive
curve for an arbitrary prescribed class $\gamma$.
\end{corollary}
\begin{proof}
Use the actual A13 curve and its integral real multiplication,
as proved in the carried A33 supplement.  A13 provides a theta
point of order $4159$ and verifies
$\langle2,-1\rangle=(\mathbf Z/4159)^\times$; A11, or the
genus-two orbit criterion KT-R2-POINT, gives the first cover.
Extend its constants to contain $\mathbf F_{128}$ and so
that the field size is one modulo $4159^b$.  The torsion
point and the needed roots of unity are then rational, so
the cyclic cover is defined after this extension.  Record
the resulting symplectic Frobenius class $\gamma$.
Apply KT-RM-FROB-AVOID with $m=4159^b$, this $\gamma$, and
$S=\{4159\}$.  The resulting second curve has the same stated
finite torsion data and no nonzero $4159$-primary theta point.
Every primitive-character doubling orbit of a degree-$4159^a$
cyclic cover would need a theta point of that exact order, so
the absence proves item~2 for every exponent.
\end{proof}

\paragraph{Scope of the new obstruction.}
For witnesses whose orders have prime factors in a prescribed
finite set, this rules out a sufficient incidence criterion
depending only on a prescribed finite congruence precision of
Frobenius and the specified real-multiplication action, even
when arbitrarily high but finite precision is permitted in
advance.  It does
not rule out a criterion involving the exact field size, a
Weil polynomial, or primes selected adaptively from the
individual curve.  It does not exhibit a curve for which every
odd cyclic cover fails.  The universal existence question of
item~7 remains unchanged.

\paragraph{Source and proof-check scope.}
The carried inputs are AN\_GRH~A19, equations (FS.3)--(FS.10),
and A37, especially its Steps~1--5 and equations
(RM.2)--(RM.4), (RM.8)--(RM.13).  Their entire relevant
proofs were read in the supplied TeX.  Their named published
moduli, descent, and Lang--Weil inputs are cited as already
accepted results, rather than newly re-proved here.  The
additional argument is the arbitrary-congruence lift (RF.3),
the valuation freezing (RF.6), and their use in the same joint
torsor specialization.
